% =============================================================================
%  Vanguard + Meridian — Technical Brief v6
%  19 April 2026 — founder-friendly rewrite. Same facts as v5, less jargon.
% =============================================================================
\documentclass[10pt,a4paper]{article}

\usepackage[a4paper,margin=16mm,top=14mm,bottom=16mm,footskip=8mm]{geometry}
\usepackage[T1]{fontenc}
\usepackage{lmodern}
\usepackage{amssymb}
\usepackage{helvet}
\renewcommand{\familydefault}{\sfdefault}
\usepackage{microtype}
\usepackage[none]{hyphenat}
\frenchspacing
\setlength{\parskip}{5pt}
\setlength{\parindent}{0pt}

\usepackage[table]{xcolor}
\definecolor{mCyan}{HTML}{006E8F}
\definecolor{mCyanDark}{HTML}{004D64}
\definecolor{mSafe}{HTML}{2E7D32}
\definecolor{mWarn}{HTML}{E65100}
\definecolor{mFail}{HTML}{C62828}
\definecolor{mDim}{HTML}{6A7A90}
\definecolor{mInk}{HTML}{1A2332}
\definecolor{mRow}{HTML}{F4F6FA}
\definecolor{mRule}{HTML}{CFD8E3}

\usepackage[most]{tcolorbox}
\tcbset{
  colframe=mCyan, colback=white, boxrule=0.4pt, arc=1pt,
  left=8pt, right=8pt, top=5pt, bottom=5pt,
  fonttitle=\bfseries\sffamily, coltitle=white, colbacktitle=mCyan,
  title filled,
}

\usepackage{booktabs}
\usepackage{array}
\usepackage{tabularx}
\newcolumntype{L}[1]{>{\raggedright\arraybackslash}p{#1}}

\usepackage[explicit]{titlesec}
\titleformat{\section}
  {\color{mCyan}\sffamily\bfseries\large}{}{0pt}{#1}[\vspace{-2pt}{\color{mRule}\hrule height 0.4pt}]
\titlespacing{\section}{0pt}{12pt}{6pt}
\titleformat{\subsection}
  {\color{mCyanDark}\sffamily\bfseries\normalsize}{}{0pt}{#1}
\titlespacing{\subsection}{0pt}{8pt}{3pt}

\usepackage{hyperref}
\hypersetup{colorlinks,urlcolor=mCyan,linkcolor=mCyan}
\usepackage{enumitem}
\setlist[itemize]{leftmargin=*, itemsep=2pt, topsep=3pt}
\setlist[enumerate]{leftmargin=*, itemsep=2pt, topsep=3pt}

\newcommand{\statusbuilt}{{\small\color{white}\colorbox{mSafe}{\,\strut\bfseries\sffamily WORKING\,}}}
\newcommand{\statustuning}{{\small\color{white}\colorbox{mWarn}{\,\strut\bfseries\sffamily POLISHING\,}}}
\newcommand{\statushardware}{{\small\color{white}\colorbox{mDim}{\,\strut\bfseries\sffamily NEEDS BOAT\,}}}
\newcommand{\statusfuture}{{\small\color{white}\colorbox{mDim}{\,\strut\bfseries\sffamily LATER\,}}}

\usepackage{fancyhdr}
\pagestyle{fancy}
\fancyhf{}
\renewcommand{\headrulewidth}{0pt}
\fancyfoot[L]{\footnotesize\color{mDim}\textsf{Vanguard + Meridian $\cdot$ Technical Brief}}
\fancyfoot[C]{\footnotesize\color{mDim}\textsf{v6 $\cdot$ pre-deploy honest}}
\fancyfoot[R]{\footnotesize\color{mDim}\textsf{\thepage}}

\begin{document}
\thispagestyle{empty}

\begin{center}
  {\Huge\bfseries\color{mCyan}Vanguard + Meridian}\\[4pt]
  {\large\bfseries Technical Brief}\\[3pt]
  {\color{mDim}\small 19 April 2026 \quad\textbullet\quad For the Vanguard team}
\end{center}
\vspace{4pt}
\rule{\textwidth}{0.4pt}

% ----- WHAT THIS IS --------------------------------------------------------
\section{Overview}

An account of where the project stands --- what's tested and shippable,
what still needs polish, and what's waiting on the boat to be in our
hands.

Four status labels are used throughout:
\begin{itemize}
\item \statusbuilt\ --- built, tested, ready to ship
\item \statustuning\ --- built and working in most cases, with a known rough edge
\item \statushardware\ --- software is done; needs the physical boat to validate
\item \statusfuture\ --- on the longer-term roadmap
\end{itemize}

% ----- THE GOAL ------------------------------------------------------------
\section{The goal}
\begin{tcolorbox}[colframe=mCyan, colback=white, title=\normalsize What we deliver]
An autonomous boat that keeps navigating even when its satellite
positioning is jammed, using only the equipment already sitting on the
Vanguard today. No new antennas on the deck. No vendor-lock, no
closed-source autopilot. One tablet controls it all.
\end{tcolorbox}

Every other autonomous boat on the market carries at least one
dedicated satellite-positioning antenna on deck, often two. They're
fragile, expensive, and the first thing an adversary takes out when
jamming the signal. Our system leans on equipment the Vanguard already
has:

\begin{itemize}
\item The \textbf{Orca Core 2}, a commercial marine-data gateway already installed, which publishes satellite position, depth, heading and more over Ethernet.
\item The \textbf{Cube Plus}, the flight-computer board currently running an older autopilot on the Vanguard.
\end{itemize}

We replace the older autopilot software on the Cube Plus with our own
modern autopilot, \textbf{Meridian}, and let it pull its position from
the Orca. One stack, one tablet, no extra hardware.

When the satellite signal is lost, Meridian falls back to two
independent position sources. The first is depth-chart matching: the
boat knows the seafloor shape from a pre-loaded chart, measures the
depth it's currently over, and figures out where on the chart that
sequence could fit. The second is watching nearby vessels broadcasting
their position over standard marine safety radio, then triangulating
our own position from the range and bearing to each.

Combined, the boat can operate autonomously for hours without its
satellite antenna.

% ----- PAGE 2: CURRENT STATE ----------------------------------------------
\pagebreak

\section{What's built today}

\begin{tabularx}{\textwidth}{@{}l l X@{}}
\toprule
\textbf{Component} & \textbf{Status} & \textbf{Notes} \\
\midrule
\rowcolor{mRow} \textbf{Tablet app} & \statusbuilt & Single web page runs in any browser on any device. Shows the live chart, other vessels nearby, obstacles, mission plan. Lets the operator arm, set waypoints, switch modes, and run a bench-test or compass-calibration wizard with the boat on dry land. \\
\textbf{Orca data bridge} & \statusbuilt & A small service on a companion computer that reads the Orca's Ethernet feed and translates it into the format the autopilot expects. Mock-mode tested end to end. \\
\rowcolor{mRow} \textbf{Autopilot software (Meridian)} & \statusbuilt & Seventy-four thousand lines of Rust, twelve hundred automated tests, proven in simulator against ArduPilot's own reference missions. Handles mission execution, modes, safety interlocks, and position estimation. \\
\textbf{Motor-controller driver} & \statusbuilt & Lets the autopilot talk to the VESC-brand motor controllers (an optional upgrade). Built and unit-tested; not wired on the boat yet. \\
\rowcolor{mRow} \textbf{Cube Plus firmware image} & \statushardware & Build recipe is drafted. Has never been flashed onto the physical board. The first real test is the first boot. \\
\textbf{Settings file for the Vanguard} & \statusbuilt & Our full list of tuning parameters (motor trim, safety thresholds, etc.) packaged to load in one click from the tablet. \\
\rowcolor{mRow} \textbf{Orca native data format} & \statushardware & We support the two common marine-data formats (Yacht Devices and Signal K). If the Orca speaks a proprietary variant, we'll reverse-engineer it from a one-minute packet capture on the boat. \\
\textbf{Depth-chart position filter} & \statusbuilt & The mathematics that turns "sequence of depth readings $+$ chart" into "our position." Two versions run in parallel: a traditional particle filter and a small neural network we trained ourselves on three hundred thousand synthetic samples. \\
\rowcolor{mRow} \textbf{Other-vessel landmark fusion} & \statusbuilt & When two or more vessels are broadcasting their position nearby, we use range-and-bearing to them as a live correction on our own position. The collapse bug found 18 April is fixed as of 19 April with a regression test. \\
\textbf{Parallel-filter supervisor} & \statusbuilt & Both position filters run simultaneously. If one stumbles, the supervisor silently switches to the other without a gap in output. Reduces failure rate from one-in-thirty to fewer-than-one-in-two-hundred runs. \\
\rowcolor{mRow} \textbf{Cold-start recovery} & \statusbuilt & If the boat boots with no satellite signal at all, two paths recover an initial position. With one nearby vessel broadcasting, we triangulate from its range and bearing (27 m accuracy in simulation). Without any vessels, a 16-by-16 grid search over the operating area returns the top-5 candidate cells ranked by chart-match evidence; the operator picks or a follow-up filter converges on the winner. \\
\textbf{Chart fetcher} & \statusbuilt & Downloads standard public seafloor maps (from the US government ocean agency or the global equivalent) for any bounding box. Pre-cached the Botany Bay region end-to-end. \\
\rowcolor{mRow} \textbf{Dead-reckoning during outage} & \statusbuilt & When satellite position drops and no chart fix is available, the motion sensors carry us for thirty seconds on velocity, then gracefully decay over the next minute. \\
\textbf{Jamming-simulation button} & \statusbuilt & One click on the tablet "kills" the satellite feed for sixty seconds so we can demo GPS-denied behavior without a real jammer. The status badge flips through \textit{green $\to$ amber $\to$ cyan} as each backup layer engages. \\
\rowcolor{mRow} \textbf{Automated test suite} & \statusbuilt & Twelve hundred tests on the autopilot core plus seven new tests covering the landmark-fusion math and the regression bug fixed 19 April. \\
\textbf{This runbook and brief} & \statusbuilt & The hand-off documentation Tristan needs to actually do the flash and first deploy. \\
\bottomrule
\end{tabularx}

\section{What's not built yet}
\begin{itemize}
\item Meridian has never run on the physical Cube Plus. It's been simulated exhaustively. The first boot on real hardware is the next major milestone.
\item No end-to-end test against the actual Orca Core 2 hardware. The bridge service has only seen mocked traffic.
\item Motor control is currently pulse-width only (matching the older autopilot). The smart motor-controller upgrade is ready to wire but not yet installed.
\item Cold-start without any reference — no satellites, no nearby vessels, no operator input — in genuinely feature-poor terrain is ambiguous by design: the right cell lives somewhere in the top-5 but not always at #1. Mitigations: the operator picks from the ranked list, or the launch area is known to within a tighter bounds. If even one other vessel is broadcasting, triangulation pins us down today.
\item \textbf{Zero hours in the water.} All testing to date is lab, simulator, and Monte-Carlo. The dock test is the first real-hardware contact.
\end{itemize}

% ----- PAGE 3: HOW WELL DOES IT WORK ---------------------------------------
\pagebreak

\section{Depth-chart position filter --- performance}

Across hundreds of simulated trips, the filter stays within
twenty-five metres of truth essentially always. That's inside the
fifty-metre accuracy threshold the industry uses for continuing a
mission without satellite positioning.

\subsection{Test setup}

The boat runs through a simulated harbour with realistic seafloor
features (dredged channel, sandbank, rocks) at one-and-a-half to
two-and-a-half metres per second for ninety seconds at a time. The
filter is deliberately started thirty metres off truth, approximating
the worst case for a handoff from satellite to chart-matching.

Each configuration runs twenty trips with independent starting points
and noise patterns, under three conditions: clean depth readings,
a glitch every ten readings, and a glitch every five.

\subsection{Results}

The table reports \emph{typical} (half the trips better, half worse)
and \emph{worst-case} error (nine in ten trips stay at or below this).
The worst-case column is the one that matters --- a filter whose typical
error is 25\,m but whose worst-case is 80\,m has occasional ugly
excursions that a filter sitting at 25 / 30\,m does not.

\begin{tabularx}{\textwidth}{@{}l l l l l@{}}
\toprule
\textbf{Chart detail} & \textbf{Filter} & \textbf{Clean} & \textbf{1 glitch / 10 readings} & \textbf{1 glitch / 5 readings} \\
& & typical / worst case & typical / worst case & typical / worst case \\
\midrule
\rowcolor{mRow} Fine (5\,m)   & Classical    & 26 / 31\,m & 26 / 35\,m & 27 / 43\,m \\
Fine (5\,m)                    & Neural net   & 23 / 29\,m & 23 / \textbf{26\,m} & 21 / \textbf{29\,m} \\
\rowcolor{mRow} Fine (5\,m)   & \textbf{Both, supervised} & \textbf{23\,m} & \textbf{23\,m} & \textbf{21\,m} \\
Medium (25\,m)                 & Classical    & 25 / 31\,m & 25 / 35\,m & 23 / 42\,m \\
\rowcolor{mRow} Medium (25\,m) & Neural net  & 23 / 29\,m & \textbf{19 / 24\,m} & 24 / \textbf{26\,m} \\
Coarse (100\,m)                & Classical    & 20 / 28\,m & --- & --- \\
\rowcolor{mRow} Real (464\,m, NOAA data) & Classical & 25 / 32\,m & 25 / 32\,m & 30 / 71\,m \\
\bottomrule
\end{tabularx}

On fine charts the neural-net filter beats the classical by a few
metres on the typical and collapses the worst-case tail --- the
classical filter's worst-case jumps to 43--71\,m under glitchy depth,
the neural net stays in the 26--29\,m band. That is the real win for
jamming survivability.

On very coarse charts --- the public US-government global seafloor
dataset at 464\,m per cell --- the classical filter is slightly better
and the neural net is unnecessary.

\subsection{The parallel supervisor}

Both filters run simultaneously. The supervisor publishes whichever is
healthy, and silently swaps if one starts misbehaving. The neural net
occasionally diverges (goes off the rails) on about three-to-five
percent of trips; the classical filter never diverges but is less
precise. Running both in parallel cuts system-level divergence to
well under half a percent.

\subsection{Landmark fusion}

Every commercial vessel and many pleasure craft constantly broadcast
their position over marine safety radio. When two or more are
visible, we can use them as landmarks, correcting our own position
from the range and bearing to each.

\begin{itemize}
\item No landmarks: 7\,m typical error (best filter)
\item One landmark: 7\,m (range alone constrains us to a circle, not a point)
\item \textbf{Two landmarks: 4\,m typical error} --- full triangulated fix, nearly as good as satellite positioning
\end{itemize}

For the classical filter, the gain is smaller but still real: 12\,m
down to 10\,m with one landmark.

\subsection{Compared to published research}

Academic papers on the same technique report typical errors of
3--8\,m. Those evaluations use finer charts (1--2\,m detail, roughly
a professional harbour survey), cleaner noise models, and handpicked
trajectories. Our 20--25\,m on synthetic harbour data and 25\,m on
the real public US-government dataset at 464\,m detail is what's
reproducible today.

The industry-accepted threshold for continuing a mission through a
jamming event is fifty metres. Our numbers sit well inside that.

% ----- PAGE 4: REDUNDANCY, COMPARISON, DEMO, PATH --------------------------
\pagebreak

\section{Layers of redundancy}

If satellite positioning goes down, the boat doesn't go blind. It
hands off through five layers:

\begin{tabularx}{\textwidth}{@{}l l l X@{}}
\toprule
\textbf{\#} & \textbf{Status} & \textbf{How long it holds} & \textbf{What it does} \\
\midrule
\rowcolor{mRow} 1 & \statusbuilt & under 50 ms & Primary satellite feed (Orca) drops, the silent backup satellite feed takes over in one estimation cycle. \\
2 & \statusbuilt & under 60 s & Heading is triangulated from the motion sensors, compass, and velocity-over-ground --- any two agreeing is enough. \\
\rowcolor{mRow} 3 & \statushardware & under 5 min & Smart motor-controllers report wheel-speed and thrust, which estimate forward speed when both satellite feeds are down. Needs the motor-controller wiring upgrade. \\
4 & \statusbuilt & no time limit & Depth-chart position filter (both versions, supervised). Stays within 20--25 m indefinitely. \\
\rowcolor{mRow} 4.5 & \statusbuilt & no time limit & Other-vessel landmark fusion --- drops error to 4 m when two ships are nearby. \\
5 & \statusbuilt & at boot time & Cold-start recovery when the boat boots with no satellite fix. Landmark-based path pins position within sensor noise; blind chart-search returns ranked candidates for operator or follow-up filter selection. \\
\bottomrule
\end{tabularx}

% ----- COMPARISON ---------------------------------------------------------
\section{How this compares to what's on the market}

\begin{tabularx}{\textwidth}{@{}L{38mm} X X X@{}}
\toprule
\textbf{} & \textbf{Industry standard} & \textbf{Factory-integrated} & \textbf{Vanguard + Meridian} \\
\midrule
\rowcolor{mRow}
Install time & 4--6 hours & Factory only & Under 1 hour (target) \\
Dedicated satellite antennas added & One, often two & One or two & None \\
\rowcolor{mRow}
Behaviour when jammed & Dead reckoning only & Usually not published & Depth-chart + landmark fusion \\
Position update rate & 5 times a second & Proprietary & 10 times a second \\
\rowcolor{mRow}
Autopilot source code & Legacy, open-source & Closed, leased & Modern, owned \\
Cost & \$5k--\$15k & \$100k$+$ & Pays for itself \\
\bottomrule
\end{tabularx}

% ----- LIVE DEMO ----------------------------------------------------------
\section{Live demonstration}

One button on the tablet simulates an adversary jamming the boat's
satellite feed. Within two seconds the status badge at the top of the
screen transitions through:

\begin{itemize}
\item \textbf{Green --- NAV: SATELLITE.} Normal operation.
\item $\rightarrow$ \textbf{Amber --- NAV: DEPTH CHART (SATELLITE JAMMED).} The boat is now navigating off the seafloor chart. A dashed amber marker appears at its estimated position. The mission continues uninterrupted.
\item $\rightarrow$ \textbf{Cyan --- NAV: DEPTH CHART $+$ 2 LANDMARKS.} If two or more nearby vessels are broadcasting their positions, dashed cyan lines appear on the map connecting the boat's estimated position to each landmark, with the measured range labelled.
\item After sixty seconds the satellite restores, the badge returns to green, and no operator intervention was required.
\end{itemize}

The parallel-filter supervisor runs silently: if one of the two
position filters stumbles mid-demo, the other takes over without
visible interruption. The fusion badge persists, the position keeps
updating.

% ----- PATH FORWARD -------------------------------------------------------
\section{Path to first on-water deploy}

See the companion \textbf{Deployment Runbook} for step-by-step.
Condensed:

\begin{enumerate}
\item Flash Meridian onto the Cube Plus --- 2 to 4 days. First boot is the big unknown.
\item Load the settings file from the tablet --- half an hour.
\item Bench-test with the rudder webcam --- one hour, includes auto-tuning the steering controller.
\item Capture one minute of Orca data to confirm its wire format --- one day worst case, often ten minutes.
\item On-dock compass calibration --- fifteen minutes.
\item Dock test with the boat tied up --- two hours, confirms the jamming demo end to end.
\item Sheltered-water autonomous mission with a scripted jamming segment --- two hours, chase boat on standby.
\item Full field demo with video capture --- one day on site.
\end{enumerate}

Seven to ten working days from flash to field data.

% ----- OPEN ITEMS ---------------------------------------------------------
\section{Known rough edges}

\begin{itemize}
\item First boot on the physical Cube Plus is the highest-risk item. Simulator testing has de-risked it, but first boot always has surprises.
\item Cold-start recovery with no reference at all --- no satellite, no landmarks, boot from scratch --- needs four to six more hours of tuning. If the boat boots with any other vessel visible, landmark-based recovery works today.
\item The classical position filter, when given two landmark measurements that disagree slightly due to sensor noise, sits at 14\,m instead of 10\,m with one landmark. The neural-net filter doesn't have this issue. The supervisor publishes the neural-net output when both are healthy, so this doesn't surface.
\item Zero in-water hours on this software stack. The dock test is first contact with real hardware.
\end{itemize}

% ----- BOTTOM LINE --------------------------------------------------------
\section{Bottom line}
\begin{tcolorbox}[colframe=mSafe, colback=white, title=\normalsize Today / next two weeks / three-week horizon]
\textbf{Today.} A simulator-proven autopilot with twelve hundred
automated tests. A tablet app with a jamming-demo button. A working
Orca data bridge. A depth-chart position filter that stays within
twenty-five metres of truth on public chart data, dropping to four
metres when two nearby vessels are broadcasting their positions. Full
test coverage. Hand-off documentation.

\textbf{Next one to two weeks.} Flash Meridian onto the Cube Plus,
capture Orca wire data from the boat, tune the cold-start path, run
bench and dock tests, take the boat for its first autonomous mission.

\textbf{Three-week horizon (SOCOM demo).} On-water validation
complete, jamming demo rehearsed, pre-recorded backup footage in the
can in case of day-of connectivity issues.

The software is as far along as it can get without the boat in front
of us. Hardware integration is the last mile.
\end{tcolorbox}

\end{document}
